%% toc: P2 21 21.4,21.5,21.6,21.7
\section{Decision procedures for the decidable fragments}
\label{sec:p02-w02-p02-035-decision-procedures}

\emph{Purpose.}  Table~\ref{tab:p02-decidable-fragments} in
\cref{sec:p02-decidable-fragments} states, row by row, the input language and
the decision complexity of four classical fragments.  This section supplies
the constructions behind all four rows: the relational monadic class,
two-variable logic \(\mathrm{FO}^2\), the Bernays--Sch\"onfinkel class and
the guarded fragment \(\mathrm{GF}\).  For each of the first three rows it
states the input, proves a finite-model bound, derives a decision procedure
from that bound, analyses its running time, and proves a matching
\(\mathsf{NEXPTIME}\) lower bound.  The guarded row needs a different method:
a tree-like model construction, a finite-model construction, a type-elimination
procedure with a doubly exponential running time, and a matching
\(2\mathsf{EXPTIME}\) lower bound.  It also records, for each row, a
nearby class whose satisfiability problem is undecidable.

\emph{Inputs and conclusions.}  The input to every procedure is a closed
first-order sentence, coded as a string over a fixed finite alphabet; its
length is written \(n\) or \(s\), following the source used in each
subsection.  The question decided is satisfiability; validity of
\(\varphi\) is the complement of satisfiability of \(\neg\varphi\).  The
complexity classes are those of
Definition~\ref{def:p02-exponential-time-classes}, and hardness is with
respect to polynomial-time many-one reductions.

\emph{Common conventions.}  Structures have nonempty domains.  A
\emph{relational} vocabulary contains relation symbols only: no constants and
no function symbols of positive arity.  Logical equality is part of the
language only where a statement says so.  The exception is
\cref{sec:p02-w02-p02-035-gf}, whose source allows individual constants; there
they are part of the input language and are removed by
\cref{thm:p02-w02-p02-035-gf-constants}, or, for the classes with a bound on
variables or arity, by the reduction behind
\cref{thm:p02-w02-p02-035-gf-bounded-constants}.  For a structure \(\mathfrak A\), a
\emph{\(k\)-type} of a tuple \((a_1,\ldots,a_k)\) is the set of atomic and
negated atomic formulas in variables \(x_1,\ldots,x_k\), equalities included
when equality is in the language, that the tuple satisfies
\cite[def.~4.1]{GradelKolaitisVardi1997}.

\subsection{The bounded torus domino problem}
\label{sec:p02-w02-p02-035-domino}

All three lower bounds below reduce from one combinatorial problem.  The
definition follows the source's torus convention: the space to be tiled wraps
around in both directions, so no border conditions are needed.

\begin{hbdefinition}{def:p02-w02-p02-035-domino}{Domino systems and the bounded torus problem}
A \emph{domino system} \(\mathcal D=(D,H,V)\) consists of a finite set \(D\)
of tiles and two relations \(H,V\subseteq D\times D\).  Write
\(\mathbb Z_s\) for the integers modulo \(s\), \(U(s,t)\defeq
\mathbb Z_s\times\mathbb Z_t\) and \(Z(t)\defeq U(t,t)\).  For a word
\(w=w_0\cdots w_{n-1}\in D^{*}\) with \(n\leq s\), \(\mathcal D\)
\emph{tiles \(U(s,t)\) with initial condition \(w\)} if there is a map
\(\tau:U(s,t)\to D\) such that, for all \((x,y)\in U(s,t)\),
\((\tau(x,y),\tau(x+1,y))\in H\), \((\tau(x,y),\tau(x,y+1))\in V\), and
\(\tau(i,0)=w_i\) for \(0\leq i<n\).  The set
\(\mathrm{DOMINO}(\mathcal D,2^n)\) consists of the words \(w\in D^{*}\) such
that \(\mathcal D\) tiles \(Z(2^{\lvert w\rvert})\) with initial condition
\(w\) \cite[defs.~6.1.1 and 6.1.5]{BorgerGradelGurevich1997}.
\index{domino system}
\end{hbdefinition}

\begin{proofboundary}{def:p02-w02-p02-035-domino-import}{Imported machine-to-tiling simulation}
The following facts about Turing machines are used as premises and are not
proved in this handbook.  (i) Every language accepted by a nondeterministic
multitape Turing machine in time \(T(n)\) is accepted by a nondeterministic
one-tape machine in time \(O(T(n)^2)\)
\cite[p.~246]{BorgerGradelGurevich1997}.  (ii) Every language accepted by a
one-tape nondeterministic machine within time \(T(n)\) and space \(S(n)\), with
\(S(n),T(n)\geq 2n\), is accepted within the same bounds by a \emph{simple}
machine in the source's sense: it never moves left of cell \(0\), and it has a
unique accepting configuration \cite[p.~243]{BorgerGradelGurevich1997}.
(iii) For every simple nondeterministic one-tape machine \(M\) with input
alphabet \(\Sigma\) there are a domino system \(\mathcal D\) and a linear-time
map \(x\mapsto w\) with \(\lvert w\rvert=\lvert x\rvert\) such that: if \(M\)
accepts \(x\) in time \(t_0\) with space \(s_0\), then \(\mathcal D\) tiles
\(U(s,t)\) with initial condition \(w\) for all \(s\geq s_0+2\) and
\(t\geq t_0+2\); and if \(M\) does not accept \(x\), then \(\mathcal D\) tiles
no \(U(s,t)\) with \(s,t\geq2\) with initial condition \(w\)
\cite[thm.~6.1.2]{BorgerGradelGurevich1997}.  The source itself gives the
tile set and horizontal relation of (iii) explicitly but leaves the full
vertical relation as an exercise \cite[pp.~244--245]{BorgerGradelGurevich1997}.
\end{proofboundary}

\begin{hblemma}{thm:p02-w02-p02-035-domino-hard}{proved=annex:p02-w02-p02-035-domino-hard}{The bounded torus domino problem is NEXPTIME-hard}
For every language \(L\in\mathsf{NEXPTIME}\) there are a domino system
\(\mathcal D\) and a polynomial-time map \(g\) such that
\(x\in L\) iff \(g(x)\in\mathrm{DOMINO}(\mathcal D,2^n)\).
\end{hblemma}

The padding in the proof makes the torus side \(2^{N}\), where \(N\) is the
length of the padded input, large enough for the whole computation.  The
lemma depends on the imported premises of
\cref{def:p02-w02-p02-035-domino-import}; the padding argument itself is
proved in full.

\subsection{The relational monadic class}
\label{sec:p02-w02-p02-035-monadic}

\begin{hbdefinition}{def:p02-w02-p02-035-monadic}{Relational monadic sentences}
A \emph{relational monadic sentence with \(m\) predicates} is a closed
first-order formula whose only nonlogical symbols are unary predicate letters
\(P_1,\ldots,P_m\).  It contains no equality, no constants and no function
symbols; quantifier prefixes and nesting are unrestricted.  The source writes
this class \([\mathit{all},(m)]\); with an unbounded number of predicates it is
\([\mathit{all},(\omega)]\), the L\"owenheim class
\cite[p.~249]{BorgerGradelGurevich1997}.  For an element \(a\) of a structure
\(\mathfrak A=(A,P_1,\ldots,P_m)\), its \emph{colour} is the word
\(c(a)=c_1\cdots c_m\in\{0,1\}^m\) with \(c_i=1\) iff
\(\mathfrak A\entails P_ia\) \cite[proof of prop.~6.2.1]{BorgerGradelGurevich1997}.
\index{monadic first-order logic!colour of an element}
\end{hbdefinition}

The colour of an element is its \(1\)-type: with no equality and no relation
of arity above one, an atomic formula in one variable \(x\) is some
\(P_ix\).

\begin{hbtheorem}{thm:p02-w02-p02-035-monadic-small-model}{proved=annex:p02-w02-p02-035-monadic-small-model}{Colour collapse and the \(2^m\) bound}
Let \(\mathfrak A\) be a structure for \(P_1,\ldots,P_m\), and let
\(\mathfrak A_{\mathrm{col}}\) be the structure whose universe is the set
\(\{c(a):a\in A\}\) of realized colours, with
\(P_i^{\mathfrak A_{\mathrm{col}}}=\{c:c_i=1\}\).  For every relational
monadic formula \(\varphi(x_1,\ldots,x_l)\) without equality and all
\(a_1,\ldots,a_l\in A\),
\[
 \mathfrak A\entails\varphi[a_1,\ldots,a_l]
 \iff
 \mathfrak A_{\mathrm{col}}\entails\varphi[c(a_1),\ldots,c(a_l)].
\]
Consequently every satisfiable relational monadic sentence with \(m\)
predicates has a model with at most \(2^m\) elements.
\end{hbtheorem}

The bound is the one the source sets as an exercise for
\([\mathit{all},(m)]\) \cite[ex.~6.2.3]{BorgerGradelGurevich1997}.  Once
equality is in the language, collapsing colours is no longer sound, because
\(\exists x\exists y\,x\neq y\) distinguishes a two-element model from its
one-element collapse.  For that case the source proves the weaker bound
\(q2^m\), where \(q\) is the quantifier rank, by keeping up to \(q\) elements
of each colour \cite[prop.~6.2.1]{BorgerGradelGurevich1997}.  The present
theorem concerns only the equality-free row of
Table~\ref{tab:p02-decidable-fragments}.

\begin{hbmethod}{def:p02-w02-p02-035-monadic-procedure}{Deciding monadic satisfiability by colour sets}
\emph{Input:} a relational monadic sentence \(\psi\) of length \(n\) with
predicates \(P_1,\ldots,P_m\).
\emph{Step 1.} Choose a nonempty set \(S\subseteq\{0,1\}^m\).
\emph{Step 2.} Form the structure \(\mathfrak B_S\) with universe \(S\) and
\(P_i^{\mathfrak B_S}=\{c\in S:c_i=1\}\).
\emph{Step 3.} Evaluate \(\psi\) in \(\mathfrak B_S\) by the recursive truth
definition.
\emph{Output:} \(\psi\) is satisfiable iff Step 3 succeeds for some choice in
Step 1.  The deterministic version tries all \(2^{2^m}-1\) choices; the
nondeterministic version guesses one.
\end{hbmethod}

\begin{hbtheorem}{thm:p02-w02-p02-035-monadic-procedure}{proved=annex:p02-w02-p02-035-monadic-procedure}{Correctness and an exponential-time bound}
The procedure of \cref{def:p02-w02-p02-035-monadic-procedure} is correct, and
its nondeterministic version runs in time \(2^{O(n^2)}\).  Hence
satisfiability of relational monadic sentences is in \(\mathsf{NEXPTIME}\),
and validity is in \(\mathsf{coNEXPTIME}\).
\end{hbtheorem}

The source obtains the sharper bound
\(\mathrm{Sat}[\mathit{all},(\omega)]_{=}\in\mathrm{NTIME}(2^{cn/\log n})\),
with equality in the language, by guessing a structure and eliminating quantifiers
one at a time over disjunctive normal forms
\cite[prop.~6.2.9 and cor.~6.2.10]{BorgerGradelGurevich1997}.  The procedure
above is cruder but suffices for membership in \(\mathsf{NEXPTIME}\).

\begin{hbdefinition}{def:p02-w02-p02-035-monadic-tiling-sentence}{The monadic tiling sentence}
Fix a domino system \(\mathcal D=(D,H,V)\) and \(w=w_0\cdots w_{n-1}\in D^n\),
\(n\geq1\).  A point \((x,y)\) of \(Z(2^n)\) is coded by the bits of
\(x=\sum_{i<n}x_i2^i\) and \(y=\sum_{i<n}y_i2^i\).  The vocabulary has unary
predicates \(X_i,X_i^{*},Y_i,Y_i^{*}\) for \(0\leq i<n\), \(N_i\) for
\(0\leq i\leq n\), and \(P_d\) for \(d\in D\), with intended meanings
\(X_iz\): \(x_i=1\); \(X_i^{*}z\): \(x_j=1\) for all \(j<i\); likewise for
\(Y_i,Y_i^{*}\); \(N_iz\): \(z=(i,0)\); and \(P_dz\): \(\tau(z)=d\).  Writing
\(\oplus\) for exclusive or, the sentence is
\[
 \psi_w\defeq\forall z\exists v\,\alpha\land\forall z\forall z'\beta,
\]
where \(\alpha\) is the conjunction of
\[
\begin{gathered}
 X_0^{*}z\land Y_0^{*}z,\qquad
 \textstyle\bigwedge_{i=1}^{n-1}\bigl(X_i^{*}z\leftrightarrow(X_{i-1}^{*}z\land X_{i-1}z)\bigr),\qquad
 \textstyle\bigwedge_{i=1}^{n-1}\bigl(Y_i^{*}z\leftrightarrow(Y_{i-1}^{*}z\land Y_{i-1}z)\bigr),\\
 \textstyle\bigwedge_{i=0}^{n-1}\bigl(X_iv\leftrightarrow(X_iz\oplus X_i^{*}z)\bigr),\qquad
 \textstyle\bigwedge_{i=0}^{n-1}\bigl(Y_iv\leftrightarrow
   [Y_iz\oplus(Y_i^{*}z\land X_{n-1}z\land X_{n-1}^{*}z)]\bigr),\\
 N_0z\leftrightarrow\textstyle\bigwedge_{i=0}^{n-1}(\neg X_iz\land\neg Y_iz),\qquad
 \textstyle\bigwedge_{i=0}^{n-1}(N_iz\leftrightarrow N_{i+1}v),
\end{gathered}
\]
and \(\beta\) is the conjunction of: exactly one \(P_dz\) with \(d\in D\);
\(H(z,z')\to\bigvee_{(d,d')\in H}(P_dz\land P_{d'}z')\);
\(V(z,z')\to\bigvee_{(d,d')\in V}(P_dz\land P_{d'}z')\); and
\(\bigwedge_{i=0}^{n-1}(N_iz\to P_{w_i}z)\).  Here
\[
 H(z,z')\defeq\textstyle\bigwedge_{i<n}(Y_iz'\leftrightarrow Y_iz)\land
   \bigwedge_{i<n}\bigl(X_iz'\leftrightarrow(X_iz\oplus X_i^{*}z)\bigr),
 \quad
 V(z,z')\defeq\textstyle\bigwedge_{i<n}(X_iz'\leftrightarrow X_iz)\land
   \bigwedge_{i<n}\bigl(Y_iz'\leftrightarrow(Y_iz\oplus Y_i^{*}z)\bigr).
\]
This is the source's construction \cite[proof of thm.~6.2.13]{BorgerGradelGurevich1997}.
\end{hbdefinition}

\begin{hbtheorem}{thm:p02-w02-p02-035-monadic-hard}{proved=annex:p02-w02-p02-035-monadic-hard}{Monadic satisfiability is NEXPTIME-complete}
For every domino system \(\mathcal D\) and \(w\in D^n\), \(\psi_w\) is
satisfiable iff \(w\in\mathrm{DOMINO}(\mathcal D,2^n)\), and \(\psi_w\) has
length \(O(n\log n)\) and is computable from \(w\) in polynomial time.
Consequently satisfiability of relational monadic sentences without equality
is \(\mathsf{NEXPTIME}\)-complete, and validity is
\(\mathsf{coNEXPTIME}\)-complete.  The lower bound already holds for sentences
of the form \(\forall x\exists y\,\alpha\land\forall x\forall y\,\beta\) with
\(\alpha,\beta\) quantifier-free and monadic.
\end{hbtheorem}

The source states the sharper lower bound
\(\mathrm{Sat}[\forall^2\exists,(\omega)]\notin\mathrm{NTIME}(2^{cn/\log n})\)
for some \(c>0\), and attributes the two-variable form to F\"urer
\cite[thm.~6.2.13 and cor.~6.2.14]{BorgerGradelGurevich1997}.  The upper and
lower bounds together match the monadic row of
Table~\ref{tab:p02-decidable-fragments}.

\begin{hbexample}{def:p02-w02-p02-035-monadic-example}{Collapsing a monadic model}
Let \(\psi\defeq\exists x(P_1x\land\neg P_2x)\land\exists x(P_2x\land\neg P_1x)
\land\forall x(P_1x\lor P_2x)\).  The structure on \(\{1,\ldots,10\}\) with
\(P_1\) the odd numbers and \(P_2\) the even numbers satisfies \(\psi\).  It
realizes exactly the colours \(10\) and \(01\), so its collapse has two
elements, \(10\) and \(01\), and satisfies \(\psi\), as
\cref{thm:p02-w02-p02-035-monadic-small-model} requires.  The bound \(2^m=4\)
is not attained here.  Now add the conjunct \(\exists x(P_1x\land P_2x)\).
The ten-element structure no longer satisfies it, but the procedure does not
inspect that structure: it searches the colour sets.  The conjunct
\(\forall x(P_1x\lor P_2x)\) rules out every set containing \(00\), and the
three existential conjuncts require \(10\), \(01\) and \(11\).  So
\(S=\{10,01,11\}\) is the only successful choice, and the extended sentence is
satisfiable by the three-element structure \(\mathfrak B_S\).
\end{hbexample}

\begin{hbwarning}{def:p02-w02-p02-035-monadic-boundary}{Where the monadic bounds stop}
The collapse fails once equality is in the language, as noted above.  Adding a
single binary predicate gives undecidability: the class of sentences
\(\psi\land\varphi\) with \(\psi\in[\forall\exists,(0,1)]\) and
\(\varphi\in[\forall^{*},(0,1)]\), whose only predicate is one binary
relation, is a conservative reduction class
\cite[thm.~3.1.23]{BorgerGradelGurevich1997}, so its satisfiability problem
is undecidable.  Adding unary function symbols keeps decidability for the
full monadic class \cite[cor.~6.2.8]{BorgerGradelGurevich1997}, but two unary
function symbols with equality already give an undecidable one-variable
class \cite[p.~11]{GradelKolaitisVardi1997}.  These boundary results are
cited, not proved here.
\end{hbwarning}

The colour-collapse proof is the finite-model side of the monadic
constituent analysis in Part~15.  There, every closed monadic formula is
shown equivalent to a disjunction of constituents \(C_S\), one for each
nonempty set \(S\) of realized profiles, and each \(C_S\) has the model with
one element per profile in \(S\)
(Theorem~\ref{thm:p15-hintikka-monadic-normal-form}).  A profile there is a
colour here, and that one-representative model is the structure
\(\mathfrak B_S\) of \cref{def:p02-w02-p02-035-monadic-procedure}.  The
constituent theorem therefore gives a second proof of the \(2^m\) bound: a
satisfiable sentence is equivalent to a nonempty disjunction of constituents,
and any one of its disjuncts \(C_S\) has a model with \(\lvert S\rvert\leq 2^m\)
elements.  It also gives a second decision procedure: compute the constituent
normal form and test whether it is nonempty.  That procedure is deterministic
and doubly exponential, since there are \(2^{2^m}-1\) constituents.

\subsection{Two-variable logic}
\label{sec:p02-w02-p02-035-fo2}

\begin{hbdefinition}{def:p02-w02-p02-035-fo2}{\(\mathrm{FO}^2\) and the Scott class}
\(\mathrm{FO}^2\) is the set of relational first-order sentences, with
equality, using at most two distinct individual variables \(x,y\), which may
be requantified \cite[sect.~2]{GradelKolaitisVardi1997}.  The size \(s\) of a
sentence is the length of its code over a fixed alphabet.  The \emph{Scott
class} consists of the sentences
\[
 \forall x\forall y\,\alpha(x,y)\ \land\
 \bigwedge_{i=1}^{m}\forall x\exists y\,\beta_i(x,y),
\]
with \(\alpha\) and every \(\beta_i\) quantifier-free; all relation symbols
have arity at most \(2\), and \(0\)-ary relation symbols are allowed
\cite[sect.~4]{GradelKolaitisVardi1997}.  An element \(a\) of a structure
\(\mathfrak A\) is a \emph{king} if no other element realizes its
\(1\)-type \cite[def.~4.2]{GradelKolaitisVardi1997}.
\index{two-variable logic!Scott class}\index{king (two-variable logic)}
\end{hbdefinition}

\begin{hbtheorem}{thm:p02-w02-p02-035-scott}{proved=annex:p02-w02-p02-035-scott}{Scott reduction}
There is a polynomial-time map sending each \(\mathrm{FO}^2\)-sentence
\(\varphi\) of size \(s\) to a sentence \(\varphi^{*}\) in the Scott class such
that: \(\varphi^{*}\) is satisfiable iff \(\varphi\) is; for every model of
\(\varphi^{*}\) there is a model of \(\varphi\) with the same universe, and
conversely; \(\varphi^{*}\) has size \(O(s\log s)\) and contains \(O(s)\)
relation symbols, all of arity at most \(2\).  Moreover, for structures with
at least two elements, each conjunct \(\forall x\exists y\,\beta_i\) may be
replaced by \(\forall x\exists y\bigl(x\neq y\land(\beta_i(x,x)\lor
\beta_i(x,y))\bigr)\) without changing truth.
\end{hbtheorem}

The source states the reduction as Scott's and records the size bound
\cite[prop.~3.1 and pp.~4--6]{GradelKolaitisVardi1997}.  The reduction is
carried out in two stages: first relation symbols of arity above two are
replaced by binary and unary ones, then every subformula receives a defining
relation symbol.

\begin{hbtheorem}{thm:p02-w02-p02-035-scott-model}{proved=annex:p02-w02-p02-035-scott-model}{Small models for the Scott class}
Let \(\theta\) be a Scott-class sentence of size \(s\) in which every
\(\beta_i\) entails \(x\neq y\), and let \(r\) be the number of relation
symbols in \(\theta\).  If \(\theta\) is satisfiable, it has a model with at
most \(3s2^r\) elements.  More precisely, if \(\mathfrak A\entails\theta\),
\(m\geq1\), \(K\) is the set of kings of \(\mathfrak A\) and \(P\) its set of
realized \(1\)-types, then \(\theta\) has a model with
\((m+1)\lvert K\rvert+3m(\lvert P\rvert-\lvert K\rvert)\) or fewer elements and
with exactly the same kings.
\end{hbtheorem}

This is the source's Theorem~4.3 \cite[thm.~4.3]{GradelKolaitisVardi1997}.
The construction keeps the kings and a \emph{royal court} of their chosen
witnesses, and replaces every other element by three banks of copies of its
\(1\)-type.  Witnesses are assigned cyclically from one bank to the next.  The
annex proof verifies both conjuncts of \(\theta\) in the new structure.

\begin{hbcorollary}{thm:p02-w02-p02-035-fo2-emp}{proved=annex:p02-w02-p02-035-fo2-emp}{Exponential model property for \(\mathrm{FO}^2\)}
There is a constant \(c\) such that every satisfiable \(\mathrm{FO}^2\)-sentence
of size \(s\) has a model with at most \(2^{cs}\) elements.  In particular
\(\mathrm{FO}^2\) has the finite-model property, and finite and unrestricted
satisfiability coincide on it.
\end{hbcorollary}

\begin{hbmethod}{def:p02-w02-p02-035-fo2-procedure}{Guess and check for \(\mathrm{FO}^2\)}
\emph{Input:} an \(\mathrm{FO}^2\)-sentence \(\varphi\) of size \(s\).
\emph{Step 1.} Compute the Scott-class sentence \(\varphi^{*}\) of
\cref{thm:p02-w02-p02-035-scott}, with \(r\) relation symbols and size
\(s^{*}\).
\emph{Step 2.} Check the finitely many one-element structures for
\(\varphi^{*}\); accept if one is a model.
\emph{Step 3.} Replace each \(\beta_i\) by the formula in the last clause of
\cref{thm:p02-w02-p02-035-scott}, obtaining \(\theta\).  Guess a structure
\(\mathfrak B\) with at most \(3s^{*}2^r\) elements, written as one truth value
for every \(0\)-ary symbol, every symbol at every element, and every binary
symbol at every pair.
\emph{Step 4.} Check \(\alpha\) at all pairs and, for each element and each
\(i\), search for a \(\beta_i\)-witness.  Accept iff all checks succeed.
\end{hbmethod}

\begin{hbtheorem}{thm:p02-w02-p02-035-fo2-complexity}{proved=annex:p02-w02-p02-035-fo2-complexity}{\(\mathrm{FO}^2\) satisfiability is NEXPTIME-complete}
For Scott-class input of size \(s\), Steps 2--4 run in nondeterministic time
\(2^{O(s/\log s)}\); for \(\mathrm{FO}^2\)-input of size \(s\), the whole
procedure runs in nondeterministic time \(2^{O(s)}\), and it is correct.
Satisfiability of \(\mathrm{FO}^2\), with or without equality, is
\(\mathsf{NEXPTIME}\)-hard, already for equality-free monadic sentences of the
form \(\forall x\exists y\,\alpha\land\forall x\forall y\,\beta\).  Hence
\(\mathrm{FO}^2\) satisfiability and finite satisfiability are
\(\mathsf{NEXPTIME}\)-complete.
\end{hbtheorem}

The upper bounds are the source's Theorem~5.3 and the completeness statement
is its Corollary~5.4 \cite[thm.~5.3 and cor.~5.4]{GradelKolaitisVardi1997}.
For the lower bound the source cites F\"urer's domino reduction to
\(\forall\forall\land\forall\exists\) sentences without proving it.  The
annex proof uses the monadic tiling sentence of
\cref{def:p02-w02-p02-035-monadic-tiling-sentence}, which has exactly that
form; the reuse of the monadic bound for two-variable logic is noted in
\cite[remark after cor.~6.2.14]{BorgerGradelGurevich1997}.

\begin{hbexample}{def:p02-w02-p02-035-king-example}{Why kings need care}
Consider
\(\forall x\exists y(Ty\land Exy)\land\forall x\exists y(Ty\land\neg Exy)\).
Without further conjuncts one can give each element two separate
\(T\)-witnesses.  If another conjunct forces the \(1\)-type containing \(Ty\)
to be realized by at most one element, that element would have to be both an
\(E\)-successor and a non-\(E\)-successor of each element.  The construction of
\cref{thm:p02-w02-p02-035-scott-model} never duplicates a king, so it can
never introduce such a conflict \cite[p.~7]{GradelKolaitisVardi1997}.
\end{hbexample}

\begin{hbwarning}{def:p02-w02-p02-035-fo2-boundary}{Where the two-variable bounds stop}
Constants may be added: they are kings, and the model bound and complexity
bound survive \cite[p.~11]{GradelKolaitisVardi1997}.  Function symbols of
positive arity may not: \(\mathrm{FO}^2\) without equality and with one unary
function symbol is undecidable \cite[p.~11]{GradelKolaitisVardi1997}.  The
three-variable fragment, even without equality, is undecidable because it
contains the \(\forall\exists\forall\) prefix class
\cite[p.~1]{GradelKolaitisVardi1997}.  The class
\(\forall\forall\forall\land\forall\exists\) contains an infinity axiom and has
an undecidable satisfiability problem \cite[p.~11]{GradelKolaitisVardi1997}.
Counting quantifiers keep decidability but lose the finite-model property
\cite[p.~11]{GradelKolaitisVardi1997}.  All of these are cited boundary
results.  Using only two banks of copies would fail: an element of one bank
and an element of the other could be each other's witnesses and require
different \(2\)-types for the same pair \cite[p.~10]{GradelKolaitisVardi1997}.
\end{hbwarning}

\subsection{The Bernays--Sch\"onfinkel class}
\label{sec:p02-w02-p02-035-bsr}

\begin{hbdefinition}{def:p02-w02-p02-035-bsr}{Bernays--Sch\"onfinkel sentences and their grounding}
A \emph{Bernays--Sch\"onfinkel sentence} is a relational prenex sentence
\[
 \varphi\defeq\exists x_1\cdots\exists x_p\forall y_1\cdots\forall y_m\,
 \theta(x_1,\ldots,x_p,y_1,\ldots,y_m)
\]
with \(\theta\) quantifier-free; the source's class \([\exists^{*}\forall^{*},
\mathit{all}]_{=}\) allows equality and the class without the subscript does
not \cite[sect.~6.2.2]{BorgerGradelGurevich1997}.  Let \(c_1,\ldots,c_p\) be
fresh constants, and if \(p=0\) let \(c_1\) be one fresh constant; write
\(q=\max(1,p)\).  The \emph{grounding} of \(\varphi\) is the set
\(G(\varphi)\) of the \(q^m\) ground instances
\(\theta(c_1,\ldots,c_p,c_{j_1},\ldots,c_{j_m})\) with
\(j_1,\ldots,j_m\in\{1,\ldots,q\}\).  An \emph{equality pattern} is an
equivalence relation \(\equiv\) on \(\{c_1,\ldots,c_q\}\).  The set
\(G_\equiv\) results from \(G(\varphi)\) by replacing every constant by the
least-indexed member of its class, and every equality atom \(c_i=c_j\) by
\(\top\) if \(c_i\equiv c_j\) and by \(\bot\) otherwise.  Without equality only
the identity pattern is used, and then \(G_\equiv=G(\varphi)\).
\index{Bernays--Schonfinkel class@Bernays--Sch\"onfinkel class!grounding}
\end{hbdefinition}

\begin{hbmethod}{def:p02-w02-p02-035-bsr-procedure}{Deciding Bernays--Sch\"onfinkel satisfiability by grounding}
\emph{Input:} \(\varphi\) as in \cref{def:p02-w02-p02-035-bsr}, of length
\(n\).
\emph{Step 1.} Write down \(G(\varphi)\) and choose an equality pattern
\(\equiv\); if \(\theta\) contains no equality, choose the identity pattern.
\emph{Step 2.} Treat each ground relation atom of \(G_\equiv\) as a
propositional variable and decide whether \(G_\equiv\) is propositionally
satisfiable.
\emph{Output:} \(\varphi\) is satisfiable iff Step 2 succeeds for some
pattern.  The nondeterministic version guesses the pattern and the valuation;
the deterministic version tries all of them.
\end{hbmethod}

\begin{hbtheorem}{thm:p02-w02-p02-035-bsr-grounding}{proved=annex:p02-w02-p02-035-bsr-grounding}{Correctness of grounding}
A Bernays--Sch\"onfinkel sentence \(\varphi\) is satisfiable iff
\(G_\equiv\) is propositionally satisfiable for some equality pattern
\(\equiv\).  A pattern with a satisfying valuation determines a model of \(\varphi\) with at
most \(\max(1,p)\) elements.
\end{hbtheorem}

The forward direction uses the small-model argument recorded in
\cref{sec:p09-presentation-n-c-005}, which is the source's
Proposition~6.2.17: the substructure induced on the existential witnesses is
again a model \cite[prop.~6.2.17]{BorgerGradelGurevich1997}.

\begin{hbtheorem}{thm:p02-w02-p02-035-bsr-complexity}{proved=annex:p02-w02-p02-035-bsr-complexity}{Bernays--Sch\"onfinkel satisfiability is NEXPTIME-complete}
The grounding procedure runs in nondeterministic time \(2^{O(n)}\).
Satisfiability of Bernays--Sch\"onfinkel sentences is
\(\mathsf{NEXPTIME}\)-hard already without equality and with prefix
\(\exists\exists\forall^{*}\).  Hence satisfiability, with or without
equality, is \(\mathsf{NEXPTIME}\)-complete, and validity of the dual
\(\forall^{*}\exists^{*}\) sentences is \(\mathsf{coNEXPTIME}\)-complete.
\end{hbtheorem}

The source proves the sharper bounds
\(\mathrm{Sat}[\exists^{*}\forall^{*},\mathit{all}]_{=}\in\mathrm{NTIME}(2^{cn})\)
and \(\mathrm{Sat}[\exists^{*}\forall^{*},\mathit{all}]\notin
\mathrm{NTIME}(2^{cn})\) for suitable constants, the second by an
\(m\)-ary encoding of the grid
\cite[thms.~6.2.19 and 6.2.21]{BorgerGradelGurevich1997}.  The annex proof of
hardness instead encodes grid points by tuples of two existentially chosen
elements.  That encoding gives polynomial-time reductions, which is what
\(\mathsf{NEXPTIME}\)-hardness requires, but not the sharper source bound.

\begin{hbexample}{def:p02-w02-p02-035-bsr-example}{A grounding with two witnesses}
Let \(\varphi\defeq\exists x_1\exists x_2\forall y\,
(Rx_1y\leftrightarrow\neg Rx_2y)\land Rx_1x_1\).  Here \(p=2\), \(m=1\), and
\(G(\varphi)\) consists of
\((Rc_1c_1\leftrightarrow\neg Rc_2c_1)\land Rc_1c_1\) and
\((Rc_1c_2\leftrightarrow\neg Rc_2c_2)\land Rc_1c_1\).  The valuation making
\(Rc_1c_1\) and \(Rc_2c_2\) true and \(Rc_2c_1\), \(Rc_1c_2\) false satisfies
both, and the structure on \(\{c_1,c_2\}\) with \(R=\{(c_1,c_1),(c_2,c_2)\}\)
is a model of \(\varphi\).
\end{hbexample}

\begin{hbwarning}{def:p02-w02-p02-035-bsr-boundary}{Where the Bernays--Sch\"onfinkel bounds stop}
The grounding is finite only because the universal block is instantiated by
constants alone.  A positive-arity function symbol produces infinitely many
ground terms, and the induced-substructure step fails.  Placing an existential
quantifier after a universal one leaves the class: already without equality,
the prefix classes \(\forall^{3}\exists\) and \(\forall\exists\forall\) of pure
predicate logic are reduction classes, and so have undecidable satisfiability
problems \cite[p.~311]{BorgerGradelGurevich1997}.  This boundary is cited, not
proved here.
\end{hbwarning}

\subsection{The guarded fragment}
\label{sec:p02-w02-p02-035-gf}

The guarded fragment lifts no restriction on arity, on the number of
variables or on the quantifier prefix.  Instead every quantifier is
relativized by an atom that contains all variables still free.  Its models
can be unravelled into tree-like models, and that is the source of both its
decidability and its finite-model property.  This subsection follows
\cite{Gradel1999}, read in its revised version of July~1998.  The
input language is that of the guarded row of
Table~\ref{tab:p02-decidable-fragments}.

\begin{hbdefinition}{def:p02-w02-p02-035-gf}{The guarded fragment \(\mathrm{GF}\)}
The vocabulary consists of relation symbols of any arity and individual
constants; there are no function symbols of positive arity, and logical
equality is available \cite[sect.~2]{Gradel1999}.  The set \(\mathrm{GF}\) of
\emph{guarded formulas} is the least set such that:
(1) every atomic formula, equalities included, is in \(\mathrm{GF}\);
(2) \(\mathrm{GF}\) is closed under \(\neg,\land,\lor,\to,\leftrightarrow\);
(3) if \(\bar x,\bar y\) are tuples of variables with \(\bar y\) nonempty,
\(\alpha(\bar x,\bar y)\) is a relation atom, \(\psi(\bar x,\bar y)\in\mathrm{GF}\),
and \(\mathrm{free}(\psi)\subseteq\mathrm{free}(\alpha)=\{\bar x,\bar y\}\),
then \(\exists\bar y(\alpha\land\psi)\) and \(\forall\bar y(\alpha\to\psi)\) are
in \(\mathrm{GF}\).  The atom \(\alpha\) is the \emph{guard}; these formulas
are written \((\exists\bar y\,.\,\alpha)\psi\) and \((\forall\bar y\,.\,\alpha)\psi\)
\cite[def.~2.1 and sect.~2]{Gradel1999}.  For \(k\geq1\),
\(\mathrm{GF}_k\) is the set of guarded formulas all of whose variables, free
or bound, are among \(x_1,\ldots,x_k\).  The length of a sentence is written
\(n\).
\index{guarded fragment}\index{guard (guarded fragment)}
\end{hbdefinition}

The source's rule~(3) asks only that the guard be atomic.  Here, as in the
catalog row, the guard is a relation atom, so an equality \(x=y\) may occur
in \(\psi\) but cannot be a guard.  The restriction is used: the first step
of the normal form below appends a fresh variable to every relation atom,
and an equality guard could not absorb it.  A nonempty \(\bar y\) loses
nothing, since for empty \(\bar y\) the formula \(\alpha\land\psi\) is already
built by rules (1) and (2).

\begin{hblemma}{thm:p02-w02-p02-035-gf-constants}{proved=annex:p02-w02-p02-035-gf-constants}{Constants can be removed}
From every guarded sentence \(\psi\) with constants \(c_1,\ldots,c_k\) one can
compute in polynomial time a constant-free guarded sentence \(\psi'\) such that,
for every set \(A\), \(\psi\) has a model with universe \(A\) iff \(\psi'\)
does.  Each \(m\)-ary symbol \(R\) of \(\psi\) is replaced by a
\((k+m)\)-ary symbol \(R^{*}\), a fresh \(k\)-ary symbol \(Z\) is added, and
\(\psi'\defeq(\exists\bar c\,.\,Z\bar c)\psi[R\bar x/R^{*}\bar c\bar x]\), the
constants \(\bar c\) now being read as variables.  The source states the
translation with a linear time bound \cite[lem.~2.3]{Gradel1999}; the annex
proves the polynomial bound used below.
\end{hblemma}

\begin{hbdefinition}{def:p02-w02-p02-035-gf-normal}{Guarded normal form}
A constant-free guarded sentence is in \emph{normal form} if it is
\[
 \psi\defeq\exists\bar x\,C\bar x\ \land\
 \bigwedge_{i\in I}(\forall\bar x\,.\,\alpha_i)\exists\bar y\,\psi_i(\bar x,\bar y),
\]
where, for each \(i\), \(\bar x=(x_1,\ldots,x_{k_i})\) with \(k_i\geq1\) and
\(\bar y=(x_{k_i+1},\ldots,x_{k_i+\ell_i})\) with \(\ell_i\geq0\) are tuples of
distinct variables, \(\alpha_i\) is a relation atom whose variables are
exactly \(x_1,\ldots,x_{k_i}\), and \(\psi_i\) is quantifier-free.  The
number of distinct variables of \(\psi\) is written \(m\)
\cite[lem.~3.1 and the remark after it]{Gradel1999}.
\end{hbdefinition}

The existential quantifier inside a conjunct need not be guarded.  The source
reaches this shape from guarded existentials \((\exists\bar y\,.\,\beta_i)\varphi_i\)
by taking the matrix \(\beta_i\land\varphi_i\), which is the step used here.

\begin{hblemma}{thm:p02-w02-p02-035-gf-normal-form}{proved=annex:p02-w02-p02-035-gf-normal-form}{Normal form}
From every constant-free guarded sentence \(\varphi\) one can compute in
polynomial time a sentence \(\psi\) in normal form such that, for every set
\(A\), \(\varphi\) has a model with universe \(A\) iff \(\psi\) does.  If
\(\varphi\) has \(v\) variables and every relation symbol of \(\varphi\) has
arity at most \(a\), then \(\psi\) has at most \(v+1\) variables and every
relation symbol of \(\psi\) has arity at most \(a+1\).
\end{hblemma}

This is the guarded part of the source's Lemma~3.1, whose proof introduces
one fresh variable and one fresh symbol \(R_\chi\) for each quantified
subformula \(\chi\) \cite[lem.~3.1]{Gradel1999}.

\begin{hbdefinition}{def:p02-w02-p02-035-gf-types}{\(\psi\)-types, witnesses and \(k\)-trees}
Let \(\psi\) be in normal form with \(m\) variables.  For \(1\leq k\leq m\),
the set \(\mathrm{At}_k(\psi)\) consists of the equalities \(x_p=x_q\)
(\(p,q\leq k\)) and every atom \(R(x_{h(z_1)},\ldots,x_{h(z_r)})\) such that
\(R(z_1,\ldots,z_r)\) is a relation atom occurring in \(\psi\) and \(h\) maps
the variables of \(\psi\) into \(\{1,\ldots,k\}\).  For a structure
\(\mathfrak A\) and \(\bar a\in A^k\), the \emph{\(\psi\)-type}
\(\mathrm{atp}^{\psi}_{\mathfrak A}(\bar a)\) is the set of formulas
\(\lambda\) and \(\neg\lambda\), \(\lambda\in\mathrm{At}_k(\psi)\), true of
\(\bar a\) in \(\mathfrak A\).  In this subsection a \emph{\(k\)-type} is any
set of this form, not the full atomic type of the common conventions.
A \(k\)-type \(s\) is a \emph{reduction} of a \(k'\)-type \(t\) if there is a
map \(\rho:\{1,\ldots,k\}\to\{1,\ldots,k'\}\) such that, whenever \(\bar a\)
realizes \(t\), \((a_{\rho(1)},\ldots,a_{\rho(k)})\) realizes \(s\).  A
\((k+\ell)\)-type \(t\) \emph{extends} a \(k\)-type \(s\) if \(s\subseteq t\),
and it extends \(s\) \emph{by \(j\) new elements} if moreover a realization of
\(t\) has \(j\) more distinct components than one of \(s\).  A \emph{witness}
for \(\psi\) is a triple \((P,t_0,(\mathrm{ext}_i)_{i\in I})\): \(P\) is a set of
\(k\)-types, \(k\leq m\), closed under reductions; \(t_0\in P\) and
\(C\bar x\in t_0\); and for each \(i\),
\(\mathrm{ext}_i:P^{(k_i)}\to P^{(k_i+\ell_i)}\) maps each \(k_i\)-type \(s\in P\)
to an extension of \(s\) containing \(\alpha_i\to\psi_i\) in the sense that
\(\alpha_i\to\psi_i\) is true under the valuation the type gives its atoms.
A structure \(\mathfrak B\) is a \emph{\(k\)-tree} if there are a tree
\(T=(V,E)\) and a map \(F\) from \(V\) to the subsets of \(B\) with at most
\(k\) elements such that (i) for every relation atom \(\alpha(x_1,\ldots,x_r)\)
and every tuple with \(\mathfrak B\entails\alpha(b_1,\ldots,b_r)\) some node
\(v\) has \(\{b_1,\ldots,b_r\}\subseteq F(v)\), and (ii) for every \(b\in B\)
the set \(\{v\in V:b\in F(v)\}\) is connected.  It has \emph{bounded degree}
if the degrees in \(T\) are bounded
\cite[defs.~3.2, 3.3 and 3.5]{Gradel1999}.
\index{guarded fragment!witness for satisfiability}\index{k-tree@\(k\)-tree}
\end{hbdefinition}

The source's types are full atomic types over the whole vocabulary.  The
\(\psi\)-types above keep only the atoms that are substitution instances of
atoms of \(\psi\).  Nothing in the criterion changes, and the count of types
becomes polynomial in the exponent once the variables or the arities are
bounded, which the \(\mathsf{EXPTIME}\) analysis needs.

\begin{hbtheorem}{thm:p02-w02-p02-035-gf-tree}{proved=annex:p02-w02-p02-035-gf-tree}{Witness criterion and tree-model property}
For a sentence \(\psi\) in normal form with \(m\) variables the following are
equivalent: (a) \(\psi\) is satisfiable; (b) \(\psi\) has a witness;
(c) \(\psi\) has a model that is an \(m\)-tree of bounded degree.
Consequently, if a constant-free guarded sentence with \(v\) variables is
satisfiable, its normal form has a model that is a \((v+1)\)-tree of bounded
degree.
\end{hbtheorem}

The proof of (b)\(\Rightarrow\)(c) is the unravelling.  Its root carries a
realization of \(t_0\).  Below a node \(v\) and for every tuple \(\bar a\)
over \(F(v)\) that contains an element new at \(v\) and satisfies a guard
\(\alpha_i\), a child \(w\) is created exactly when
\(\mathrm{ext}_i(\mathrm{atp}(\bar a))\) needs new elements; \(F(w)\) is
\(\bar a\) together with fresh elements realizing that extension.  Every
relation atom on a tuple that lies in no single set \(F(v)\) is made false.
The source proves sufficiency for the loosely guarded fragment
\cite[lem.~3.4, lem.~3.6 and thm.~3.7]{Gradel1999}; the annex proof is
written for the guarded fragment and for \(\psi\)-types.

\begin{proofboundary}{def:p02-w02-p02-035-herwig-import}{Imported extension theorem}
The following theorem is used as a premise and is not proved in this
handbook.  Call a tuple \(\bar a\) of a structure \emph{live} if it satisfies
some relation atom \(\beta(x_1,\ldots,x_r)\) in which all of
\(x_1,\ldots,x_r\) occur.  For every finite structure \(\mathfrak A\) with
finite relational vocabulary there is a finite structure \(\mathfrak A^{+}\)
containing \(\mathfrak A\) as an induced substructure such that (i) every
partial isomorphism of \(\mathfrak A\) extends to an automorphism of
\(\mathfrak A^{+}\), and (ii) for every live tuple \(\bar a\) of
\(\mathfrak A^{+}\) some automorphism \(g\) of \(\mathfrak A^{+}\) maps
\(\bar a\) into \(A\).  The source states this as Herwig's theorem and refers
to Herwig's papers for its proof \cite[def.~3.8 and thm.~3.9]{Gradel1999}.
\end{proofboundary}

\begin{hbtheorem}{thm:p02-w02-p02-035-gf-fmp}{proved=annex:p02-w02-p02-035-gf-fmp}{Finite-model property}
Every satisfiable guarded sentence, with or without constants and equality,
has a finite model.  Hence finite and unrestricted satisfiability coincide on
\(\mathrm{GF}\).
\end{hbtheorem}

The construction cuts the tree model of \cref{thm:p02-w02-p02-035-gf-tree} to
a finite subtree \(S\) that is closed under predecessors and contains a node
of every isomorphism type of the bounded structures \(F(v)\).  It then takes
the extension \(\mathfrak A^{+}\) of the part of the tree model on \(S\).  A
guard satisfied in \(\mathfrak A^{+}\) is moved into that part by an
automorphism.  Its witness obligation is met in a node of \(S\) isomorphic to
the node where the tree model met it, and the composed automorphism carries
the witness back \cite[thm.~3.10]{Gradel1999}.  The source leaves open
whether the loosely guarded fragment has the finite-model property
\cite[remark after thm.~3.10]{Gradel1999}.

\begin{hbmethod}{def:p02-w02-p02-035-gf-procedure}{Deciding guarded satisfiability by type elimination}
\emph{Input:} a guarded sentence \(\varphi\) of length \(n\).
\emph{Step 1.} Remove constants by \cref{thm:p02-w02-p02-035-gf-constants}
and bring the result into normal form \(\psi\) by
\cref{thm:p02-w02-p02-035-gf-normal-form}; let \(m\) be its number of
variables.
\emph{Step 2.} List the sets \(\mathrm{At}_k(\psi)\), \(k\leq m\), and the set
\(P_0\) of all \(k\)-types, \(k\leq m\).  A set \(t\) that contains exactly one
of \(\lambda,\neg\lambda\) for each \(\lambda\in\mathrm{At}_k(\psi)\) is a
\(k\)-type iff its equalities form an equivalence relation \(\approx\) on
\(\{1,\ldots,k\}\) and any two atoms of \(\mathrm{At}_k(\psi)\) with the same
symbol and \(\approx\)-equivalent arguments are both in \(t\) or both negated
in \(t\).
\emph{Step 3.} Given \(P_j\), let \(P_{j+1}\) consist of the \(t\in P_j\) such
that every reduction of \(t\) with at most \(m\) variables lies in \(P_j\),
and, for every \(i\) with \(t\) a \(k_i\)-type and \(\alpha_i(x_1,\ldots,x_{k_i})\in t\),
some \((k_i+\ell_i)\)-type \(t'\in P_j\) extends \(t\) and makes \(\psi_i\)
true.  Stop when
\(P_{j+1}=P_j\), and call this set \(P^{*}\).
\emph{Output:} \(\varphi\) is satisfiable iff some \(t_0\in P^{*}\) contains
\(C\bar x\).
\end{hbmethod}

\begin{hbtheorem}{thm:p02-w02-p02-035-gf-upper}{proved=annex:p02-w02-p02-035-gf-upper}{Correctness and running time of type elimination}
The procedure of \cref{def:p02-w02-p02-035-gf-procedure} is correct and runs
in deterministic time \(2^{2^{O(n\log n)}}\), so guarded satisfiability is in
\(2\mathsf{EXPTIME}\).  For every fixed \(k\), on constant-free guarded
sentences with at most \(k\) variables, and on constant-free guarded
sentences whose relation symbols have arity at most \(k\), it runs in time
\(2^{n^{O(1)}}\), so these problems are in \(\mathsf{EXPTIME}\).
\end{hbtheorem}

The source obtains the same upper bounds from an alternating test that
guesses types and universally chooses reductions within exponential space,
then uses the equality of alternating exponential space and deterministic
doubly exponential time \cite[thm.~4.3]{Gradel1999}.  The procedure above
computes the largest set of types that satisfies the witness conditions
directly, so no alternation theorem is needed for the upper bound.

The bounded statements of the source, Corollary~4.6, concern its language
with individual constants \cite[cor.~4.6]{Gradel1999}.  The translation of
\cref{thm:p02-w02-p02-035-gf-constants} does not preserve those bounds: it
adds every constant as an argument of every relation symbol and as a variable,
so both the arity and the number of variables grow by the number of constants.
The bounded classes with constants need a different reduction.

\begin{hbcorollary}{thm:p02-w02-p02-035-gf-bounded-constants}{proved=annex:p02-w02-p02-035-gf-bounded-constants}{Bounded guarded sentences with constants}
For every fixed \(k\geq1\), satisfiability of guarded sentences with
individual constants and at most \(k\) variables, and of guarded sentences
with individual constants whose relation symbols have arity at most \(k\), is
decidable in deterministic time \(2^{n^{O(1)}}\).  Both problems are therefore
in \(\mathsf{EXPTIME}\).
\end{hbcorollary}

The reduction splits the universe of a model into the elements named by
constants and the rest.  Each variable is either sent to a named element or
kept as a variable over the unnamed elements.  An atom \(R(t_1,\ldots,t_r)\)
becomes an atom of a new symbol \(R_\kappa\), where \(\kappa\) records which
argument places carry which constant, applied to the remaining variables
only.  Arity and variable count do not grow.  The number of cases for one
quantified block is polynomial when the number of variables in a guard is
bounded.  To keep the nesting of quantifiers from multiplying these cases,
the split is applied after the sentence has been flattened as in
\cref{thm:p02-w02-p02-035-gf-normal-form}.  The annex proof also uses a
property of guarded sentences that holds with or without constants: adding
an element that lies in no relation does not change their truth.

\begin{proofboundary}{def:p02-w02-p02-035-alternation-import}{Imported facts on alternating machines}
The lower bounds use the following premises, which are not proved in this
handbook.  (i) For every space bound \(s(n)\geq\log n\), the problems
decided by alternating Turing machines in space \(s(n)\) are exactly those
decided deterministically in time \(2^{d\,s(n)}\) for some \(d\); in
particular alternating exponential space equals \(2\mathsf{EXPTIME}\) and
alternating polynomial space equals \(\mathsf{EXPTIME}\)
\cite[defs.~4.1 and 4.2 and sect.~4.1]{Gradel1999}.  (ii) Every problem
decided by an alternating machine within space \(S(n)\) is decided within the
same space by a one-tape alternating machine \(M\) whose states are
existential, universal, accepting or rejecting, every computation of which
halts, and every non-final configuration of which has exactly two successor
configurations \cite[proof of thm.~4.4]{Gradel1999}.  The machine model has
tape cells numbered from \(0\), and the head of \(M\) stays in cells
\(0,\ldots,S(n)-1\).
\end{proofboundary}

\begin{hbdefinition}{def:p02-w02-p02-035-gf-accept}{The guarded acceptance sentence}
Let \(M\) be as in premise (ii) with space \(2^{s}\), \(s=p(n)\) for a
polynomial \(p\), state set \(Q\) with initial state \(q_0\), tape alphabet
\(\Gamma\) with blank \(\square\), and successor transitions
\(\delta_1(q,\sigma),\delta_2(q,\sigma)\in Q\times\Gamma\times\{-1,0,1\}\) for
non-final \(q\).  Fix an input \(w=w_0\cdots w_{n-1}\).  Nodes of a
computation tree are coded by pairs \((a,b)\); a cell number \(r<2^{s}\) is
coded by the tuple \(\hat r(a,b)\in\{a,b\}^{s}\) whose \(i\)-th entry is \(b\)
iff bit \(i-1\) of \(r\) is \(1\).  A unary \(N\) marks \(a\) and not \(b\).
With \(\bar x=(x_1,\ldots,x_s)\) and \(\bar x'\) similar, put
\[
\begin{aligned}
 \mathit{succ}(\bar x,\bar x')&\defeq\textstyle\bigvee_{i}\bigl(
   \bigwedge_{j<i}(\neg Nx_j\land Nx'_j)\land Nx_i\land\neg Nx'_i\land
   \bigwedge_{j>i}(Nx_j\leftrightarrow Nx'_j)\bigr),\\
 \mathit{diff}(\bar x,\bar x')&\defeq\textstyle\bigvee_{i}\neg(Nx_i\leftrightarrow Nx'_i).
\end{aligned}
\]
The vocabulary has \(N\), binary \(T\) (node), \(R\) (root) and \(Q_q\)
(\(q\in Q\)); \((s+2)\)-ary \(U\), \(Z\) (head) and \(P_\sigma\)
(\(\sigma\in\Gamma\)); \(4\)-ary \(S_1,S_2\) (successors); and \((4s+4)\)-ary
\(E\).  Let \(\mathrm{EXIST}(a,b)\), \(\mathrm{UNIV}(a,b)\) and
\(\mathrm{REJECT}(a,b)\) be the disjunctions of \(Q_qab\) over existential,
universal and rejecting \(q\).  \(\mathrm{NEXT}_i(a,b,c,d)\) is the
conjunction of
\[
 (\forall\bar x\bar x'\bar y\bar y'\,.\,E\bar x\bar x'ab\bar y\bar y'cd)
 \bigl((Z\bar x'ab\land\mathit{diff}(\bar x,\bar x'))\to
 \textstyle\bigwedge_{\sigma}(P_\sigma\bar xab\to P_\sigma\bar ycd)\bigr)
\]
and, for each non-final \(q\) and each \(\sigma\) with
\(\delta_i(q,\sigma)=(q',\sigma',\mu)\), the formula
\[
 (\forall\bar x\bar x'\bar y\bar y'\,.\,E\bar x\bar x'ab\bar y\bar y'cd)
 \bigl((Q_qab\land Z\bar xab\land P_\sigma\bar xab\land\gamma_\mu)\to
 (Q_{q'}cd\land P_{\sigma'}\bar ycd\land Z\bar y_\mu cd)\bigr),
\]
where
\(\gamma_1=\mathit{succ}(\bar x,\bar x')\), \(\gamma_{-1}=\mathit{succ}(\bar x',\bar x)\),
\(\gamma_0=\top\), \(\bar y_{\pm1}=\bar y'\) and \(\bar y_0=\bar y\).  Then
\(\mathrm{ACCEPT}_{M,w}\) is the conjunction of the following nine
sentences, numbered for reference:
\begin{enumerate}
\item \((\exists ab\,.\,Rab)Tab\);
\item \((\forall ab\,.\,Tab)\bigl(Na\land\neg Nb\land U(a,\ldots,a,a,b)\bigr)\);
\item \((\forall\bar xab\,.\,U\bar xab)\bigwedge_{i}U\bar x[x_i/b]ab\);
\item the sentence
  \[
   (\forall ab\,.\,Tab)\bigl((\mathrm{EXIST}(a,b)\lor\mathrm{UNIV}(a,b))\to
   \theta_1\land\theta_2\bigr),
  \]
  where \(\theta_i\defeq(\exists cd\,.\,S_iabcd)Tcd\);
\item for \(i=1,2\), the sentence
  \((\forall abcd\,.\,S_iabcd)\bigl(E\bar aab\bar ccd\land\epsilon\bigr)\), where
  \[
   \epsilon\defeq(\forall\bar u\bar v\,.\,E\bar uab\bar vcd)
   \textstyle\bigwedge_{j\leq2s}E\bar u[u_j/b]ab\bar v[v_j/d]cd;
  \]
\item the sentence
  \(
   (\forall ab\,.\,Rab)\bigl(Q_{q_0}ab\land Z\hat0ab\land
   \bigwedge_{r<n}P_{w_r}\hat rab\land\iota\bigr)
  \), where
  \[
   \iota\defeq(\forall\bar x\,.\,U\bar xab)
   \bigl(\textstyle\bigwedge_{r<n}\mathit{diff}(\hat r,\bar x)\to P_\square\bar xab\bigr);
  \]
\item \((\forall abcd\,.\,S_1abcd)\bigl(\mathrm{EXIST}(a,b)\to
  \mathrm{NEXT}_1(a,b,c,d)\lor\mathrm{NEXT}_2(a,b,c,d)\bigr)\);
\item for \(i=1,2\), \((\forall abcd\,.\,S_iabcd)
  \bigl(\mathrm{UNIV}(a,b)\to\mathrm{NEXT}_i(a,b,c,d)\bigr)\);
\item \((\forall ab\,.\,Tab)\neg\mathrm{REJECT}(a,b)\).
\end{enumerate}
where \(\bar a\) and \(\bar c\) repeat \(a\) and \(c\) \(2s\) times and
\(\bar u,\bar v\) are \(2s\)-tuples of variables.  The symbols \(T,R,U,Z,
P_\sigma,Q_q,S_i,E,N\), the formula \(\mathit{succ}\),
and the names \(\mathrm{EXIST}\), \(\mathrm{UNIV}\), \(\mathrm{REJECT}\),
\(\mathrm{NEXT}_i\) and \(\mathrm{ACCEPT}_{M,w}\) follow the source
\cite[proof of thm.~4.4]{Gradel1999}; \(\mathit{diff}\) replaces the source's
use of \(\mathit{less}\) and of tuple inequality.
\end{hbdefinition}

\begin{hbtheorem}{thm:p02-w02-p02-035-gf-hard}{proved=annex:p02-w02-p02-035-gf-hard}{Guarded satisfiability is 2EXPTIME-complete}
\(\mathrm{ACCEPT}_{M,w}\) is a guarded sentence without equality and without
constants, computable from \(w\) in polynomial time, and it is satisfiable
iff \(M\) accepts \(w\).  Hence satisfiability of guarded sentences, with or
without equality, is \(2\mathsf{EXPTIME}\)-complete.  Moreover, for every
\(k\geq2\), satisfiability of guarded sentences with at most \(k\) variables,
and of guarded sentences whose relation symbols have arity at most \(k\), is
\(\mathsf{EXPTIME}\)-hard already for equality-free sentences in
\(\mathrm{GF}_2\) with relation symbols of arity at most \(2\).  With or
without individual constants, these bounded problems are
\(\mathsf{EXPTIME}\)-complete.
\end{hbtheorem}

The source's sentence differs in two places.  It forces in every model a
binary tree of depth \(2^{2^{n}}\) through an \(n\)-ary depth counter \(D\)
and requires two successors at every node.  As printed, every successor
must carry a counter one larger, an increment needs a zero bit, and the
counter is bounded; an infinite chain of successors is therefore impossible,
so the printed conjunction has no model.  The sentence above
drops the counter.  It requires successors only at existential and universal
configurations, and it relies on premise (ii): every computation halts.  The
source also uses equality in \(\mathit{less}\) and \(\mathit{succ}\); the
predicate \(N\) is the source's own device for removing equality, applied
here from the start \cite[proof of thm.~4.4 and cor.~4.5]{Gradel1999}.  For the
\(\mathsf{EXPTIME}\) lower bound the source sketches an encoding of
linear-space configurations by monadic predicates
\cite[cor.~4.6]{Gradel1999}; the annex writes that encoding out in
\(\mathrm{GF}_2\).

\begin{hbexample}{def:p02-w02-p02-035-gf-example}{An infinite tree model and a finite model}
Let \(\varphi\defeq(\exists xy\,.\,Exy)\top\land(\forall xy\,.\,Exy)(\exists z\,.\,Eyz)\top\),
where \(\top\) abbreviates \(Exy\lor\neg Exy\) inside the first guard and
\(Eyz\lor\neg Eyz\) inside the second.  Every \(E\)-edge must be followed by
another, so the natural numbers with \(E=\{(r,r+1)\}\) form a model.  It is a
\(2\)-tree of bounded degree: the path \(v_0,v_1,\ldots\) with
\(F(v_r)=\{r,r+1\}\) satisfies conditions (i) and (ii).  Up to the auxiliary symbols of the normal form, this is the shape the
unravelling of \cref{thm:p02-w02-p02-035-gf-tree} produces: each node below
the root holds one old element and one new witness.  All structures
\(F(v_r)\) are isomorphic, so the finite subtree \(S=\{v_0\}\) already
represents every isomorphism type.  The finite-model construction then takes
an extension of the edge \(\{0,1\}\) in which the partial isomorphism
\(0\mapsto1\) extends to an automorphism; such an extension needs a further
element to receive the image of \(1\).  The directed \(E\)-cycle on
\(\{0,1,2\}\) is one such structure: rotations extend each partial
isomorphism and move each edge onto \((0,1)\).  It is a finite model of
\(\varphi\), and it also satisfies the guarded conjuncts
\((\forall xy\,.\,Exy)\neg Eyx\) and \((\forall x\,.\,Exx)\neg\top\), with \(\top\) read as \(Exx\lor\neg Exx\), which rule out
two-element and one-element cycles.
\end{hbexample}

\begin{hbwarning}{def:p02-w02-p02-035-gf-boundary}{Where the guarded bounds stop}
A conjunction of atoms is not a guard: transitivity,
\(\forall xyz(Exy\land Eyz\to Exz)\), is neither guarded nor loosely guarded,
since no single atom contains \(x\) and \(z\)
\cite[sect.~2]{Gradel1999}.  Small extensions are undecidable.  Sentences
\(\psi\land\mathrm{functional}[D]\) with \(\psi\in\mathrm{GF}_3\) without
equality, where \(\mathrm{functional}[D]\) says that the binary relation \(D\)
is a partial function, form a conservative reduction class
\cite[thm.~5.8]{Gradel1999}; so does \(\mathrm{GF}_3\) with counting
quantifiers \cite[cor.~5.9]{Gradel1999}, and so do the sentences
\(\psi\land\mathrm{transitive}[E_0,\ldots,E_m]\) with \(\psi\in\mathrm{GF}_3\)
\cite[thm.~5.10]{Gradel1999}.  Their satisfiability problems are therefore
undecidable.  These boundary results are cited, not proved here.  The
loosely guarded fragment shares the tree-model property and the
\(2\mathsf{EXPTIME}\) bound, but its finite-model property was left open by
the source \cite[remark after thm.~3.10]{Gradel1999}.
\end{hbwarning}

The guarded fragment contains the first-order translation of propositional modal
logic; that translation and the bisimulation side of guarded logic are in
\cref{sec:p09-presentation-n-c-007} and
Theorem~\ref{thm:p09-modal-guarded}.

\subsection{Relationships}
\label{sec:p02-w02-p02-035-relationships}

The monadic, two-variable and Bernays--Sch\"onfinkel rows share one method: a satisfiable sentence has a model whose
size is bounded by a computable function of the sentence, a nondeterministic
machine guesses such a model, and a model check confirms it.  The bounds
differ in kind.  For monadic sentences the bound \(2^m\) counts colours; for
\(\mathrm{FO}^2\) the bound \(3s2^r\) counts \(1\)-types times witness
obligations; for Bernays--Sch\"onfinkel sentences the bound \(\max(1,p)\)
counts existential witnesses, so the exponential cost lies in the number of
ground instances rather than in the model.  The guarded row does not fit this
pattern.  Its finite models come from an extension theorem, not from a
computable size bound proved here, so the procedure searches sets of types
instead of models.  The extra exponential is the number of types: a type
records atoms of unbounded arity over up to \(m\) variables.  Bounding the
variables or the arity removes it, which is why the bounded guarded classes
fall to \(\mathsf{EXPTIME}\).  The finite-model consequences of
these bounds bear on the distinction between finite and unrestricted validity
in \cref{sec:p02-w02-p02-036-finite-validity}.  The decision and complexity
summary in Part~14 lists the same rows in Table~\ref{tab:p14-dc-finite}.
